\documentclass[border=4pt]{standalone}
\usepackage{amsmath,amssymb,mathtools,xcolor,varwidth}
\definecolor{fanpurple}{HTML}{8E24AA}
\definecolor{centerorange}{HTML}{E8770A}
\begin{document}
\begin{varwidth}{28em}
\large\bfseries
Proposition I proved that a central force produces \textcolor{fanpurple}{equal triangles SAB, SBC, SCD — all equal areas} swept in equal times. Proposition II is its converse: now the equal areas are the PREMISE, and centripetal force is the conclusion. About the fixed \textcolor{centerorange}{centre S} the body follows the polygonal path $A\,B\,C\,D\,E$, and we are GIVEN that it describes areas $\propto$ times. This is the observational direction — exactly what Kepler measured for Mars: if you OBSERVE $\triangle SAB=\triangle SBC=\triangle SCD$, must the force point to $S$?
\end{varwidth}
\end{document}
